**Conventions.**
All graphs are **finite, simple, undirected** graphs on a **labelled** vertex set (V(G)={1,\dots,n}).
A **cycle** means a **simple cycle** (no repeated vertices) of length (\ell\ge 3).

**Cycle set.** For a graph (G) on (n) vertices define
[
S(G);:=;{\ell\in{3,4,\dots,n}:\text{ (G) contains a cycle of length (\ell)}}.
]
A set (A\subseteq{3,\dots,n}) is called a **cycle set on (n) vertices** if there exists an (n)-vertex graph (G) with (S(G)=A).

**Counting function.** Let
[
f(n);:=;\bigl|{S(G):; G \text{ is a graph on }n\text{ vertices}}\bigr|
]
(the number of distinct cycle sets realized by (n)-vertex graphs).

**Claims to prove.**

1. (f(n)=o(2^n)) as (n\to\infty), i.e.
   [
   \lim_{n\to\infty}\frac{f(n)}{2^n}=0.
   ]
2. (\displaystyle \lim_{n\to\infty}\frac{f(n)}{2^{n/2}}=\infty), i.e.
   [
   \forall M>0;\exists N;\forall n\ge N:\quad f(n)>M\cdot 2^{n/2}.
   ]

**Edge cases.** For (n\le 2), ({3,\dots,n}=\emptyset), so the only cycle set is (\emptyset) and (f(1)=f(2)=1). For asymptotics we only care about (n\to\infty).

---
